Un beurre est rance si le pourcentage en masse d'acide butanoïque qu'il contient
est supérieur à $4\%$, c'est-à-dire qu'il y a plus de $\qty{4}{\gram}$ d'acide
butanoïque dans $\qty{100}{\gram}$ de beurre.

\begin{donnees}
    \begin{itemize}
        \item $M(\ce{C3H7CO2H})=\qty{88}{\gram\per\mol}$.
    \end{itemize}
\end{donnees}

Dans un bécher, on introduit $m_{b}=\qty{10,0}{\gram}$ de beurre fondu auquel on
ajoute de l'eau distillée. On agite afin de dissoudre dans l'eau la totalité de
l'acide butanoïque $\ce{C3H7CO2H}$ présent dans le beurre. On obtient une
solution aqueuse $(S)$ d'acide butanoïque de concentration molaire $C$ et de
volume $V_{0}=\qty{1,0}{\liter}$.

On dose le volume $V=\qty{10,0}{\mL}$ de la solution $(S)$ par une solution
aqueuse d'hydroxyde de sodium $\ce{Na^+_{(aq)} + HO^-_{(aq)}}$ de concentration
molaire $C_{B}=\qty{4,0e-3}{\M}$.
          
\begin{minipage}{0.60\linewidth}
    Le suivi pH-métrique du dosage permet d'obtenir la courbe
    ${\mathrm{pH}=f(V_{B})}$ (Figure ci-contre).  On considère que seul l'acide
    butanoïque réagit avec le réactif titrant.
    \begin{enumerate}
        \item Écrire l'équation de la réaction du dosage sachant
              qu'elle est totale.
        \item Déterminer graphiquement la valeur du volume $V_{B,\mathrm{E}}$ à
              l'équivalence.
        \item Calculer la valeur de $C$.
        \item Déterminer la masse de l'acide butanoïque présent dans
              la masse $m_{b}=\qty{10,0}{\gram}$ du beurre. Le beurre
              étudié est-il rance~? Justifier.
    \end{enumerate}
\end{minipage}
\hfill
\begin{minipage}{0.38\linewidth}
    \flushright
    \begin{figure}[H]
        \centering
        \includegraphics[width=\textwidth]{2026-03-23_122336-}
    \end{figure}
\end{minipage}

